\SourceLeaf{008}
\begin{center}\small
DEMONSTRATIO NOVA THEOREMATIS\\[5pt]
\textbf{OMNEM FUNCTIONEM ALGEBRAICAM RATIONALEM INTEGRAM}\\[5pt]
UNIUS VARIABILIS\\[5pt]
IN FACTORES REALES PRIMI VEL SECUNDI GRADUS RESOLVI POSSE.\\[8pt]
\rule{1.5cm}{0.3pt}
\end{center}
\Article{1}
Quaelibet aequatio algebraica determinata reduci potest ad formam
\[
x^m+Ax^{m-1}+Bx^{m-2}+\text{etc.}+M=0
\]
ita ut $m$ sit numerus integer positivus. Si partem primam huius aequationis per
$X$ denotamus, aequationique $X=0$ per plures valores inaequales ipsius $x$ satisfieri supponimus, puta ponendo $x=\alpha$, $x=\beta$, $x=\gamma$ etc. functio $X$ per
productum e factoribus $x-\alpha$, $x-\beta$, $x-\gamma$ etc. divisibilis erit. Vice versa, si
productum e pluribus factoribus simplicibus $x-\alpha$, $x-\beta$, $x-\gamma$ etc. functionem
$X$ metitur: aequationi $X=0$ satisfiet, aequando ipsam $x$ cuicunque quantitatum $\alpha$, $\beta$, $\gamma$ etc. Denique si $X$ producto ex $m$ factoribus talibus simplicibus
aequalis est (sive omnes diversi sint, sive quidam ex ipsis identici): alii factores
simplices praeter hos functionem $X$ metiri non poterunt. Quamobrem aequatio
$m^{\mathrm{ti}}$ gradus plures quam $m$ radices habere nequit; simul vero patet, aequationem
$m^{\mathrm{ti}}$ gradus \textit{pauciores} radices habere posse, etsi $X$ in $m$ factores simplices resolubilis sit: si enim inter hos factores aliqui sunt identici, multitudo modorum diversorum aequationi satisfaciendi necessario minor erit quam $m$. Attamen concinnitatis caussa geometrae dicere maluerunt, aequationem in hoc quoque casu $m$
radices habere, et tantummodo quasdam ex ipsis aequales inter se evadere: quod
utique sibi permittere poterunt.
\SourceLeaf{009}
\Article{2}
Quae hucusque sunt enarrata, in libris algebraicis sufficienter demonstrantur neque rigorem geometricum uspiam offendunt. Sed nimis praepropere et sine
praevia demonstratione solida adoptavisse videntur analystae theorema cui tota
fere doctrina aequationum superstructa est: \textit{Quamvis functionem talem ut $X$ semper in $m$ factores simplices resolvi posse}, sive hoc quod cum illo prorsus conspirat,
\textit{quamvis aequationem $m^{\mathrm{ti}}$ gradus revera habere $m$ radices}. Quum iam in aequationibus secundi gradus saepissime ad tales casus perveniatur, qui theoremati huic
repugnant: algebraistae, ut hos illi subiicerent, coacti fuerunt, fingere quantitatem quandam imaginariam cuius quadratum sit $-1$, et tum agnoverunt, si quantitates formae $a+b\sqrt{-1}$ perinde concedantur ut reales, theorema non modo pro
aequationibus secundi gradus verum esse, sed etiam pro cubicis et biquadraticis.
Hinc vero neutiquam inferre licuit, admissis quantitatibus formae $a+b\sqrt{-1}$ cuivis aequationi quinti superiorisve gradus satisfieri posse, aut uti plerumque exprimitur (quamquam phrasin lubricam minus probarem) radices cuiusvis aequationis ad formam $a+b\sqrt{-1}$ reduci posse. Hoc theorema ab eo, quod in titulo
huius scripti enunciatum est, nihil differt. si ad rem ipsam spectas, huiusque demonstrationem novam rigorosam tradere. constituit propositum praesentis dissertationis.

Ceterum ex eo tempore, quo analystae comperti sunt, infinite multas aequationes esse, quae nullam omnino radicem haberent, nisi quantitates formae
$a+b\sqrt{-1}$ admittantur, tales quantitates fictitiae tamquam peculiare quantitatum genus, quas \textit{imaginarias} dixerunt, ut a \textit{realibus} distinguerentur, consideratae
et in totam analysin introductae sunt; quonam iure? hoc loco non disputo. --- Demonstrationem meam absque omni quantitatum imaginariarum subsidio absolvam, etsi eadem libertate, qua omnes recentiores analystae usi sunt, etiam mihi
uti liceret.

\Article{3}
Quamvis ea, quae in plerisque libris elementaribus tamquam demonstratio
theorematis nostri afferuntur, tam levia sint, tantumque a rigore geometrico abhorreant, ut vix mentione sint digna: tamen, ne quid desideretur, paucis
illa attingam. `Ut demonstrent, quamvis aequationem
\[
x^m+Ax^{m-1}+Bx^{m-2}+\text{etc.}+M=0
\]
\SourceLeaf{010}
sive $X=0$, revera habere $m$ radices, suscipiunt probare, $X$ in $m$ factores
simplices resolvi posse. Ad hunc finem assumunt $m$ factores simplices $x-\alpha$,
$x-\beta$, $x-\gamma$ etc. ubi $\alpha$, $\beta$, $\gamma$ etc. adhuc sunt incognitae, productumque ex illis
aequale ponunt functioni $X$. Tum ex comparatione coëfficientium deducunt $m$
aequationes, ex quibus incognitas $\alpha$, $\beta$, $\gamma$ etc. determinari posse aiunt, quippe
quarum multitudo etiam sit $m$. Scilicet $m-1$ incognitas eliminari posse, unde
emergere aequationem, quae, quam placuerit, incognitam solam contineat.' Ut
de reliquis, quae in tali argumentatione reprehendi possent, taceam, quaeram
tantummodo, unde certi esse possimus, ultimam aequationem revera ullam radicem habere? Quidni fieri posset, ut neque huic ultimae aequationi neque propositae, ulla magnitudo in toto quantitatum realium atque imaginariarum ambitu
satisfaciat? --- Ceterum periti facile perspicient, hanc ultimam aequationem necessario cum proposita \textit{omnino identicam} fore, siquidem calculus rite fuerit institutus; scilicet eliminatis incognitis $\beta$, $\gamma$ etc. aequationem
\[
\alpha^m+A\alpha^{m-1}+B\alpha^{m-2}+\text{etc.}+M=0
\]
prodire debere. Plura de isto ratiocinio exponere necesse non est.

Quidam auctores, qui debilitatem huius methodi percepisse videntur, tamquam \textit{axioma} assumunt, quamvis aequationem revera habere radices, si non possibiles, impossibiles. Quid sub quantitatibus possibilibus et impossibilibus intelligi velint, haud satis distincte exposuisse videntur. Si quantitates possibiles
idem denotare debent ut reales, impossibiles idem ut imaginariae: axioma illud
neutiquam admitti potest, sed necessario demonstratione opus habet. Attamen
in illo sensu expressiones accipiendae non videntur, sed axiomatis mens haec potius videtur esse: `Quamquam nondum sumus certi, necessario dari $m$ quantitates reales vel imaginarias, quae alicui aequationi datae $m^{\mathrm{ti}}$ gradus satisfaciant,
tamen aliquantisper hoc supponemus; nam si forte contingeret, ut tot quantitates reales et imaginariae inveniri nequeant, certe effugium patebit, ut dicamus
reliquas esse impossibiles.' Si quis hac phrasi uti mavult quam simpliciter dicere, aequationem in hoc casu tot radices non habituram, a me nihil obstat: at
si tum his radicibus impossibilibus ita utitur tamquam aliquid veri sint, et e. g.
dicit, summam omnium radicum aequationis $x^m+Ax^{m-1}+\text{etc.}=0$ esse
$=-A$, etiamsi impossibiles inter illas sint (quae expressio proprie significat,
\textit{etiamsi aliquae deficiant}): hoc neutiquam probare possum. Nam radices impossibiles\SourceLeaf{011}, in tali sensu acceptae, tamen sunt radices, et tum axioma illud nullo modo
sine demonstratione admitti potest, neque inepte dubitares, annon aequationes
exstare possint, quae ne impossibiles quidem radices habeant?\HistoricalNote{1}

\Article{4}
Antequam aliorum geometrarum demonstrationes theorematis nostri recenseam, et quae in singulis reprehendenda mihi videantur, exponam: observo sufficere si tantummodo ostendatur, omni aequationi quantivis gradus
\[
x^m+Ax^{m-1}+Bx^{m-2}+\text{etc.}+M=0
\]
sive $X=0$ (ubi coëfficientes $A,B$ etc. reales esse supponuntur) ad minimum
uno modo satisfieri posse per valorem ipsius $x$ sub forma $a+b\sqrt{-1}$ contentum. Constat enim, $X$ tunc divisibilem fore per factorem realem secundi gradus
$xx-2ax+aa+bb$, si $b$ non fuerit $=0$, et per factorem realem simplicem
$x-a$, si $b=0$. In utroque casu quotiens erit realis, et inferioris gradus quam
$X$; et quum hic eadem ratione factorem realem primi secundive gradus habere debeat\SourceLeaf{012}, patet, per continuationem huius operationis functionem $X$ tandem in factores reales simplices vel duplices resolutum iri, aut, si pro singulis factoribus
realibus duplicibus binos imaginarios simplices adhibere mavis, in $m$ factores
simplices.

\Article{5}
Prima theorematis demonstratio illustri geometrae \textsc{d'Alembert} debetur, \textit{Recherches sur le calcul intégral, Histoire de l'Acad. de Berlin, Année} 1746.
p. 182 sqq. Eadem extat in \textsc{Bougainville}, \textit{Traité du calcul intégral, à Paris} 1754.
p. 47 sqq. Methodi huius praecipua momenta haec sunt.

Primo ostendit, si functio quaecunque $X$ quantitatis variabilis $x$ fiat $=0$
aut pro $x=0$ aut pro $x=\infty$, atque valorem infinite parvum realem positivum nancisci possit tribuendo ipsi $x$ valorem realem: hanc functionem etiam valorem infinite parvum realem negativum obtinere posse per valorem ipsius $x$ vel
realem vel sub forma imaginaria $p+q\sqrt{-1}$ contentum. Scilicet designante $\Omega$
valorem infinite parvum ipsius $X$, et $\omega$ valorem respondentem ipsius $x$, asserit
$\omega$ per seriem valde convergentem $a\Omega^\alpha+b\Omega^\beta+c\Omega^\gamma$ etc. exprimi posse, ubi exponentes $\alpha$, $\beta$, $\gamma$ etc. sint quantitates rationales continuo crescentes, et quae adeo
ad minimum in distantia certa ab initio positivae evadant, terminosque, in quibus adsint, infinite parvos reddant. Iam si inter omnes hos exponentes nullus
occurrat, qui sit fractio denominatoris paris, omnes terminos seriei reales fieri tum
pro positivo tum pro negativo valore ipsius $\Omega$; si vero quaedam fractiones denominatoris paris inter illos exponentes reperiantur, constare, pro valore negativo
ipsius $\Omega$ terminos respondentes in forma $p+q\sqrt{-1}$ contentos esse. Sed propter infinitam seriei convergentiam in casu priori sufficere, si terminus primus (i. e.
maximus) solus retineatur, in posteriori ultra eum terminum, qui partem imaginariam primus producat, progredi opus non esse.

Per similia ratiocinia ostendi posse, si $X$ valorem negativum infinite parvum ex valore reali ipsius $x$ assequi possit: functionem illam valorem realem
positivum infinite parvum ex valore reali ipsius $x$ vel ex imaginario sub forma
$p+q\sqrt{-1}$ contento adipisci posse.

Hinc secundo concludit, etiam valorem aliquem realem finitum ipsius $X$
dari, in casu priori negativum, in posteriori positivum, qui ex valore imaginario
ipsius $x$ sub forma $p+q\sqrt{-1}$ contento produci possit.

\SourceLeaf{013}
Hinc sequitur, si $X$ sit talis functio ipsius $x$, quae valorem realem $V$
ex valore ipsius $x$ reali $v$ obtineat, atque etiam valorem realem quantitate infinite parva vel maiorem vel minorem ex valore reali ipsius $x$ assequatur, eandem
etiam valorem realem quantitate infinite parva atque adeo finita vel minorem vel
maiorem quam $V$ (resp.) recipere posse, tribuendo ipsi $x$ valorem sub forma
$p+q\sqrt{-1}$ contentum. Hoc nullo negotio ex praecc. derivatur, si pro $X$ substitui concipitur $V+Y$, et pro $x$, $v+y$.

Tandem affirmat ill. \textsc{d'Alembert}, si $X$ totum intervallum aliquod inter duos
valores reales $R,S$ percurrere posse supponatur (i. e. tum ipsi $R$, tum ipsi $S$,
tum omnibus valoribus realibus intermediis aequalis fieri), tribuendo ipsi $x$ valores semper in forma $p+q\sqrt{-1}$ contentos; functionem $X$ quavis quantitate finita reali adhuc augeri vel diminui posse (prout $S>R$ vel $S<R$), manente $x$
semper sub forma $p+q\sqrt{-1}$. Si enim quantitas realis $U$ detur (inter quam
et $R$ supponitur $S$ iacere), cui $X$ per talem valorem ipsius $x$ aequalis fieri non
posset, necessario valorem \textit{maximum} ipsius $X$ dari (scilicet quando $S>R$; minimum vero, quando $S<R$), puta $T$, quem ex valore ipsius $x$, $p+q\sqrt{-1}$,
consequeretur, ita ut ipsi $x$ nullus valor sub simili forma contentus tribui posset,
qui functionem $X$ vel minimo excessu propius versus $U$ promoveret. Iam si in
hac aequatione inter $X$ et $x$ pro $x$ ubique substituatur $p+q\sqrt{-1}$, atque tum pars
realis, tum pars, quae factorem $\sqrt{-1}$ implicet, hoc omisso, cifrae aequentur:
ex duabus aequationibus hinc prodeuntibus (in quibus $p,q$ et $X$ cum constantibus permixtae occurrent) per eliminationem duas alias elici posse, in quarum altera $p,X$ et constantes reperiantur, altera a $p$ libera solas $q,X$ et constantes
involvat. Quamobrem quum $X$ per valores reales ipsarum $p,q$ omnes valores ab
ipsa $R$ usque ad $T$ percurrerit, per praecc. $X$ versus valorem $U$ adhuc propius accedere posse tribuendo ipsis $p,q$ valores tales $\alpha+\gamma\sqrt{-1}$, $\beta+\delta\sqrt{-1}$ resp. Hinc
vero fieri $x=\alpha-\delta+(\gamma+\beta)\sqrt{-1}$, i. e. adhuc sub forma $p+q\sqrt{-1}$ esse, contra hyp.

Iam si $X$ functionem talem ut $x^m+Ax^{m-1}+Bx^{m-2}+\text{etc.}+M$ denotare supponitur, nullo negotio perspicitur, ipsi $x$ tales valores reales tribui posse,
ut $X$ totum aliquod intervallum inter duos valores reales percurrat. Quare $x$
valorem aliquem sub forma $p+q\sqrt{-1}$ contentum talem nancisci poterit,
unde $X$ fiat $=0$. Q. E. D.\HistoricalNote{2}
\SourceLeaf{014}
\Article{6}
Quae contra demonstrationem \textsc{d'Alembertianam} obiici posse videntur, ad
haec fere redeunt.

1. Ill. d'A. nullum dubium movet de \textit{existentia} valorum ipsius $x$, quibus valores dati ipsius $X$ respondeant, sed illam supponit, solamque \textit{formam} istorum
valorum investigat.

Quamvis vero haec obiectio per se gravissima sit, tamen hic ad solam dictionis formam pertinet, quae facile ita corrigi potest, ut illa penitus destruatur.

2. Assertio, $\omega$ per talem seriem qualem ponit semper exprimi posse, certo
est falsa, si $X$ etiam functionem quamlibet transcendentem designare debet (uti
d'A. pluribus locis innuit). Hoc e. g. manifestum est, si ponitur $X=e^{\frac1x}$, sive
$x=\frac1{\log X}$. Attamen si demonstrationem ad eum casum restringimus, ubi $X$ est
functio algebraica ipsius $x$ (quod in praesenti negotio sufficit), propositio utique
est vera. --- Ceterum d'A. nihil pro confirmatione suppositionis suae attulit; cel.
\textsc{Bougainville} supponit, $X$ esse functionem algebraicam ipsius $x$, et ad inventionem seriei parallelogrammum \textsc{Newtonianum} commendat.

3. Quantitatibus infinite parvis liberius utitur, quam cum geometrico rigore consistere potest aut saltem nostra aetate (ubi illae merito male audiunt) ab
analysta scrupuloso concederetur, neque etiam saltum a valore infinite parvo ipsius $\Omega$ ad finitum satis luculenter explicat. Propositionem suam, $\Omega$ etiam valorem aliquem finitum consequi posse, non tam ex possibilitate valoris infinite
parvi ipsius $\Omega$ concludere videtur quam inde potius, quod denotante $\Omega$ quantitatem valde parvam, propter magnam seriei convergentiam, quo plures termini
seriei accipiantur, eo propius ad valorem verum ipsius $\omega$ accedatur, aut,
quo plurium partium summa pro $\omega$ accipiatur, eo exactius aequationi, quae relationem inter $\omega$ et $\Omega$ sive $x$ et $X$ exhibeat, satisfactum iri. Praeterea quod tota
haec argumentatio nimis vaga videtur, quam ut ulla conclusio rigorosa inde colligi possit: observo, utique dari series, quae quantumvis parvus valor quantitati,
\SourceLeaf{015}
secundum cuius potestates progrediuntur, tribuatur, nihilominus semper divergant, ita ut si modo satis longe continuentur, ad terminos quavis quantitate data
maiores pervenire possis\HistoricalNote{3}. Hoc evenit, quando coëfficientes seriei progressionem
hypergeometricam constituunt. Quamobrem necessario demonstrari debuisset, talem seriem hypergeometricam in casu praesenti provenire non posse.

Ceterum mihi videtur, ill. d'A. hic non recte ad series infinitas confugisse,
hasque ad stabiliendum theorema hoc fundamentale doctrinae aequationum haud
idoneas esse.

4. Ex suppositione, $X$ obtinere posse valorem $S$ neque vero valorem $U$,
nondum sequitur, inter $S$ et $U$ necessario valorem $T$ iacere, quem $X$ attingere
sed non superare possit. Superest adhuc alius casus: scilicet fieri posset, ut inter $S$ et $U$ limes situs sit, ad quem accedere quidem quam prope velis possit $X$,
ipsum vero nihilominus numquam attingere. Ex argumentis ab ill. d'A. allatis
tantummodo sequitur, $X$ omnem valorem, quem attigerit, adhuc quantitate
finita superare posse, puta quando evaserit $=S$, adhuc quantitate aliqua finita
$\Omega$ augeri posse; quo facto, novum incrementum $\Omega'$ accedere, tunc iterum augmentum $\Omega''$ etc., ita ut quotcunque incrementa iam adiecta sint, nullum pro ultimo
haberi debeat, sed semper aliquod novum accedere possit. At quamvis \textit{multitudo}
incrementorum possibilium nullis limitibus sit circumscripta: tamen utique fieri
posset, ut si incrementa $\Omega,\Omega',\Omega''$ etc. continuo decrescerent, nihilominus summa
$S+\Omega+\Omega'+\Omega''$ etc. limitem aliquem numquam attingeret, quotcunque termini
considerentur.

Quamquam hic casus occurrere non potest, quando $X$ designat functionem
algebraicam integram ipsius $x$: tamen sine demonstratione, hoc fieri non posse,
methodus necessario pro incompleta habenda est. Quando vero $X$ est functio
transcendens, sive etiam algebraica fracta, casus ille utique locum habere potest,
e. g. semper quando valori cuidam ipsius $X$ valor infinite magnus ipsius $x$ respondet\SourceLeaf{016}. Tum methodus \textsc{d'Alembertiana} non sine multis ambagibus, et in quibusdam
casibus nullo forsan modo, ad principia indubitata reduci posse videtur.

Propter has rationes demonstrationem \textsc{d'Alembertianam} pro satisfaciente
habere nequeo. Attamen hoc non obstante verus demonstrationis nervus probandi
per omnes obiectiones neutiquam infringi mihi videtur, credoque eidem fundamento (quamvis longe diversa ratione, et saltem maiori circumspicientia) non solum demonstrationem rigorosam theorematis nostri superstrui, sed ibinde omnia
peti posse, quae circa aequationum \textit{transscendentium} theoriam desiderari queant.
De qua re gravissima alia occasione fusius agam; conf. interim infra art. 24.

\Article{7}
Post \textsc{d'Alembertum} ill. \textsc{Euler} disquisitiones suas de eodem argumento promulgavit, \textit{Recherches sur les racines imaginaires des équations, Hist. de l'Acad.
de Berlin A.} 1749, p. 223 sqq. Methodum duplicem hic tradidit: prioris summa
continetur in sequentibus.

Primo ill. E. suscipit demonstrare, si $m$ denotet quamcunque dignitatem
numeri 2, functionem $x^{2m}+Bx^{2m-2}+Cx^{2m-3}+\text{etc.}+M=X$ (in qua coëfficiens termini secundi est $=0$) semper in duos factores reales resolvi posse, in
quibus $x$ usque ad $m$ dimensiones ascendat. Ad hunc finem duos factores assumit,
\[
x^m-ux^{m-1}+\alpha x^{m-2}+\beta x^{m-3}+\text{etc.},\quad\text{et}\quad x^m+ux^{m-1}+\lambda x^{m-2}+\mu x^{m-3}\ \text{etc.}
\]
ubi coëfficientes $u,\alpha,\beta$ etc. $\lambda,\mu$ etc. adhuc incogniti sunt, horumque productum
aequale ponit functioni $X$. Tum coëfficientium comparatio suppeditat $2m-1$
aequationes, manifestoque demonstrari tantummodo debet, incognitis $u,\alpha,\beta$ etc.
$\lambda,\mu$ etc. (quarum multitudo etiam est $2m-1$) tales valores reales tribui posse, qui
aequationibus illis satisfaciant. Iam E. affirmat, si primo $u$ tamquam cognita
consideretur, ita ut multitudo incognitarum unitate minor sit quam multitudo
aequationum, his secundum methodos algebraicas notas rite combinatis omnes
$\alpha,\beta$ etc. $\lambda,\mu$ etc. rationaliter et sine ulla radicum extractione per $u$ et coëfficientes $B,C$ etc. determinari posse, adeoque valores reales nancisci, simulac $u$
realis fiat. Praeterea vero omnes $\alpha,\beta$ etc. $\lambda,\mu$ etc. eliminari poterunt, ita ut
prodeat aequatio $U=0$, ubi $U$ erit functio integra solius $u$ et coëfficientium
cognitorum. Hanc aequationem, ipsam per methodum eliminationis vulgarem evolvere, opus immensum foret, quando aequatio proposita $X=0$ est gradus aliquantum\SourceLeaf{017} alti; et pro gradu indeterminato, plane impossibile (iudice ipso E. p. 239).
Attamen hic sufficit, unam illius aequationis proprietatem novisse, scilicet quod
terminus ultimus in $U$ (qui incognitam $u$ non implicat) necessario est negativus,
unde sequi constat, aequationem ad minimum unam radicem realem habere, sive
$u$ et proin etiam $\alpha,\beta$ etc. $\lambda,\mu$ etc. ad minimum uno modo realiter determinari
posse: illam vero proprietatem per sequentes reflexiones confirmare licet. Quum
$x^m-ux^{m-1}+\alpha x^{m-2}+\text{etc.}$ supponatur esse factor functionis $X$: necessario $u$
erit summa $m$ radicum aequationis $X=0$, adeoque totidem valores habere debebit, quot modis diversis ex $2m$ radicibus $m$ excerpi possunt, sive per principia calculi combinationum $\frac{2m\mathbin{.}\overline{2m-1}\mathbin{.}\overline{2m-2}.\,.\,.\,.\,.\overline{m+1}}{1\mathbin{.}2\mathbin{.}3\ldots m}$ valores. Hic numerus semper erit impariter par (demonstrationem haud difficilem supprimo): si itaque ponitur $=2k$, ipsius semissis $k$ impar erit; aequatio $U=0$ vero erit gradus
$2k^{\mathrm{ti}}$. Iam quoniam in aequatione $X=0$ terminus secundus deest: summa
omnium $2m$ radicum erit 0; unde patet, si summa quarumcunque $m$ radicum
fuerit $+p$, reliquarum summam fore $-p$, i. e. si $+p$ est inter valores ipsius
$u$, etiam $-p$ inter eosdem erit. Hinc E. concludit, $U$ esse productum ex $k$
factoribus duplicibus talibus $uu-pp$, $uu-qq$, $uu-rr$ etc., denotantibus $+p$,
$-p$, $+q$, $-q$ etc. omnes $2k$ radices aequationis $U=0$, unde, propter multitudinem imparem horum factorum, terminus ultimus in $U$ erit quadratum producti $pqr$ etc. signo negativo affectum. Productum autem $pqr$ etc. semper ex
coëfficientibus $B,C$ etc. rationaliter determinari potest, adeoque necessario erit
quantitas realis. Huius itaque quadratum signo negativo affectum certo erit quantitas negativa. Q. E. D.

Quum hi duo factores reales ipsius $X$ sint gradus $m^{\mathrm{ti}}$ atque $m$ potestas
numeri 2: eadem ratione uterque rursus in duos factores reales $\frac12m$ dimensionum
resolvi poterit. Quoniam vero per repetitam dimidiationem numeri $m$ necessario
tandem ad binarium pervenitur, manifestum est, per continuationem operationis
functionem $X$ tandem in factores reales secundi gradus resolutam haberi.

Quodsi vero functio talis proponitur, in qua terminus secundus non deest,
puta $x^{2m}+Ax^{2m-1}+Bx^{2m-2}+\text{etc.}+M$, designante etiamnum $2m$ potestatem
binariam, haec per substitutionem $x=y-\frac{A}{2m}$ transibit in similem functionem
termino secundo carentem. Unde facile concluditur, etiam illam functionem in
factores reales secundi gradus resolubilem esse.

Denique proposita functione gradus $n^{\mathrm{ti}}$, designante $n$ numerum, qui non
\SourceLeaf{018}
est potestas binaria: ponatur potestas binaria proxime maior quam $n$, $=2m$,
multipliceturque functio proposita per $2m-n$ factores simplices reales quoscunque.
Ex resolubilitate producti in factores reales secundi gradus, nullo negotio
derivatur, etiam functionem propositam in factores reales secundi vel primi gradus
resolubilem esse debere.

\Article{8}
Contra hanc demonstrationem obiici potest

1. Regulam, secundum quam E. concludit, ex $2m-1$ aequationibus
$2m-2$ incognitas $\alpha,\beta$ etc. $\lambda,\mu$ etc. omnes rationaliter determinari posse,
neutiquam esse generalem, sed saepissime exceptionem pati. Si quis e. g. in art. 3.,
aliqua incognitarum tamquam cognita spectata, reliquas per hanc et coëfficientes
datos rationaliter exprimere tentat, facile inveniet, hoc esse impossibile, nullamque
quantitatem incognitam aliter quam per aequationem $m-1^{\mathrm{ti}}$ gradus
determinari posse. Quamquam vero hic statim a priori perspici potest, illud
necessario ita evenire debuisse: tamen merito dubitari posset, annon etiam in casu
praesenti pro quibusdam valoribus $m$ res eodem modo se habeat, ut incognitae
$\alpha,\beta$ etc. $\lambda,\mu$ etc. ex $u,B,C$ etc. aliter quam per aequationem gradus forsan
maioris quam $2m$ determinari nequeant. Pro eo casu, ubi aequatio $X=0$ est
quarti gradus, E. valores rationales coëfficientium per $u$ et coëfficientes datos eruit;
verum vero etiam in omnibus aequationibus altioribus fieri posse, utique explicatione
ampliori egebat. --- Ceterum operae pretium esse videtur, in formulas illas,
quae $\alpha,\beta$ etc. rationaliter per $u,B,C$ etc. exprimant, profundius et
generalissime inquirere; de qua re aliisque ad eliminationis theoriam (argumentum
haudquaquam exhaustum) pertinentibus alia occasione fusius agere suscipiam.

2. Etiamsi autem demonstratum fuerit, cuiusvis gradus sit aequatio $X=0$,
semper formulas inveniri posse, quae ipsas $\alpha,\beta$ etc. $\lambda,\mu$ etc. rationaliter per
$u,B,C$ etc. exhibeant: tamen certum est, pro valoribus quibusdam determinatis
coëfficientium $B,C$ etc. formulas illas \textit{indeterminatas} evadere posse, ita ut non
solum impossibile sit, incognitas illas rationaliter ex $u,B,C$ etc. definire, sed
adeo revera quibusdam in casibus valori alicui reali ipsius $u$ nulli valores reales
ipsarum $\alpha,\beta$ etc. $\lambda,\mu$ etc. respondeant. Ad confirmationem huius rei brevitatis
gratia ablego lectorem ad diss. ipsam E., ubi p. 236 aequatio quarti gradus fusius
explicata est. Statim quisque videbit, formulas pro coëfficientibus $\alpha,\beta$ indeterminatas\SourceLeaf{019} fieri, si $C=0$ et pro $u$ assumatur valor $0$, illorumque valores non solum
sine extractione radicum assignari non posse, sed adeo ne reales quidem esse,
si fuerit $BB-4D$ quantitas negativa. Quamquam vero in hoc casu $u$ adhuc
alios valores reales habere, quibus valores reales ipsarum $\alpha,\beta$ respondeant,
facile perspici potest: tamen vereri aliquis posset, ne huius difficultatis enodatio
(quam E. omnino non attigit) in aequationibus altioribus multo maiorem operam
facessat. Certe haec res in demonstratione exacta neutiquam silentio praeteriri
debet.

3. Ill. E. supponit tacite, aequationem $X=0$ habere $2m$ radices, harumque
summam statuit $=0$, ideo quod terminus secundus in $X$ abest. Quomodo
de hac licentia (qua omnes auctores de hoc argumento utuntur) sentiam, iam
supra art. 3 declaravi. Propositio, summam omnium radicum aequationis alicuius
coëfficienti primo, mutato signo, aequalem esse, ad alias aequationes applicanda
non videtur, nisi quae radices habent: iam quum per hanc ipsam demonstrationem
evinci debeat, aequationem $X=0$ revera radices habere, haud permissum
videtur, harum existentiam supponere. Sine dubio ii, qui huius paralogismi
fallaciam nondum penetraverunt, respondebunt, \textit{hic non demonstrari, aequationi
$X=0$ satisfieri posse} (nam hoc dicere vult expressio, eam habere radices),
\textit{sed tantummodo, ipsi per valores ipsius $x$ sub forma $a+b\sqrt{-1}$ contentos satisfieri
posse; illud vero tamquam axioma supponi.} At quum aliae quantitatum formae,
praeter realem et imaginariam $a+b\sqrt{-1}$ concipi nequeant, non satis luculentum
videtur, quomodo id, quod demonstrari debet, ab eo, quod tamquam axioma
supponitur, differat; quin adeo si possibile esset adhuc alias formas quantitatum
excogitare, puta formam $F,F',F''$ etc.: tamen sine demonstratione admitti non
deberet, cuius aequationi per aliquem valorem ipsius $x$ aut realem, aut sub forma
$a+b\sqrt{-1}$, aut sub forma $F$, aut sub $F'$ etc. contentum satisfieri posse. Quamobrem
axioma illud alium sensum habere nequit quam hunc: Cuivis aequationi
satisfieri potest \textit{aut} per valorem realem incognitae, \textit{aut} per valorem imaginarium
sub forma $a+b\sqrt{-1}$ contentum, \textit{aut} forsan per valorem sub forma alia hucusque
ignota contentum, \textit{aut} per valorem, qui sub nulla omnino forma continetur.
Sed quomodo huiusmodi quantitates, de quibus ne ideam quidem fingere potes
--- vera umbrae umbra --- summari aut multiplicari possint, hoc ea perspicuitate,
quae in mathesi semper postulatur, certo non intelligitur\HistoricalNote{4}.

\SourceLeaf{020}
Ceterum conclusiones, quas E. ex suppositione sua elicuit, per has obiectiones
haudquaquam suspectas reddere volo; quin potius certus sum, illas per methodum
neque difficilem neque ab \textsc{Euleriana} multum diversam ita comprobari
posse, ut nemini vel minimus scrupulus superesse debeat. Solam \textit{formam} reprehendo,
quae quamvis in \textit{inveniendis} novis veritatibus magnae utilitatis esse possit,
tamen in \textit{demonstrando}, coram publico, minime probanda videtur.

4. Pro demonstratione assertionis, productum $pqr$ etc. ex coëfficientibus
in $X$ \textit{rationaliter} determinari posse, ill. E. nihil omnino attulit. Omnia, quae
hac de re in aequationibus \textit{quarti gradus} explicat, haec sunt (ubi $a,b,c,d$ sunt
radices aequationis propositae $x^4+Bxx+Cx+D=0$):

\lq On m'objectera sans doute, que j'ai supposé ici, que la quantité $pqr$ était
une quantité réelle, et que son quarré $ppqqrr$ était affirmatif; ce qui était encore
douteux, vu que les racines $a,b,c,d$ étant imaginaires, il pourrait bien arriver,
que le quarré de la quantité $pqr$, qui en est composée, fut négatif. Or je réponds
à cela que ce cas ne saurait jamais avoir lieu; car quelque imaginaires que
soient les racines $a,b,c,d$, on sait pourtant, qu'il doit y avoir $a+b+c+d=0$:
$ab+ac+ad+bc+bd+cd=B$; $abc+abd+acd+bcd=-C$\HistoricalNote{5}; $abcd=D$,
ces quantités $B,C,D$ étant réelles. Mais puisque $p=a+b$, $q=a+c$,
$r=a+d$, leur produit $pqr=(a+b)(a+c)(a+d)$ est déterminable \textit{comme on
sait}, par les quantités $B,C,D$, et sera par conséquent réel, tout comme nous
avons vu, qu'il est effectivement $pqr=-C$, et $ppqqrr=CC$. On reconnaîtra
aisément de même, que dans les plus hautes équations cette même circonstance
doit avoir lieu, et qu'on ne saurait me faire des objections de ce
côté.\rq
Conditionem, productum $pqr$ etc. \textit{rationaliter} per $B,C$ etc. determinari
posse, E. nullibi adiecit, attamen semper subintellexisse videtur, quum absque
illa demonstratio nullam vim habere possit. Iam verum quidem est in aequationibus
quarti gradus, si productum $(a+b)(a+c)(a+d)$ evolvatur, obtineri
$aa(a+b+c+d)+abc+abd+acd+bcd=-C$, attamen non satis perspicuum
videtur, quomodo in omnibus aequationibus superioribus productum rationaliter
\SourceLeaf{021}
per coëfficientes determinari possit. Clar. \textsc{de Foncenex}, qui primus hoc observavit
(\textit{Miscell. phil. math. soc. Taurin.} T. I. p. 117), recte contendit, sine demonstratione
rigorosa huius propositionis methodum omnem vim perdere, illam vero
satis difficilem sibi videri confitetur, et quam viam frustra tentaverit, enarrat\HistoricalNote{6}.
Attamen haec res haud difficulter per methodum sequentem (cuius summam
addigitare tantummodo hic possum) absolvitur: Quamquam in aequationibus quarti
gradus non satis clarum est, productum $(a+b)(a+c)(a+d)$ per coëfficientes
$B,C,D$ determinabile esse, tamen facile perspici potest, idem productum
etiam esse $=(b+a)(b+c)(b+d)$, nec non $=(c+a)(c+b)(c+d)$,
denique etiam $=(d+a)(d+b)(d+c)$. Quare productum $pqr$ erit quadrans summae
$(a+b)(a+c)(a+d)+(b+a)(b+c)(b+d)+(c+a)(c+b)(c+d)+(d+a)(d+b)(d+c)$,
quam, si evolvatur, fore functionem rationalem integram radicum $a,b,c,d$ talem,
in quam omnes eadem ratione ingrediantur, nullo negotio a priori praevideri
potest. Tales vero functiones semper rationaliter per coëfficientes aequationis,
cuius radices sunt $a,b,c,d$, exprimi possunt. --- Idem etiam manifestum est,
si productum $pqr$ sub hanc formam redigatur:
\[
\tfrac12(a+b-c-d)\times\tfrac12(a+c-b-d)\times\tfrac12(a+d-b-c)
\]
quod productum evolutum omnes $a,b,c,d$ eodem modo implicaturum esse facile
praevideri potest. Simul periti facile hinc colligent, quomodo hoc ad altiores
aequationes applicari debeat. --- Completam demonstrationis expositionem, quam
hic apponere brevitas non permittit, una cum uberiori disquisitione de functionibus
plures variabiles eodem modo involventibus ad aliam occasionem mihi reservo.

Ceterum observo, praeter has quatuor obiectiones, adhuc quaedam alia in
demonstratione E. reprehendi posse, quae tamen silentio praetereo, ne forte censor
nimis severus esse videar, praesertim quum praecedentia satis ostendere videantur,
demonstrationem in ea quidem forma, in qua ab E. proposita est, pro
completa neutiquam haberi posse.

Post hanc demonstrationem, E. adhuc aliam viam theorema pro aequationibus,
quarum gradus non est potestas binaria, ad talium aequationum resolutionem
reducendi ostendit: attamen quum methodus haec pro aequationibus quarum
